\section{A limitation of cycle, theta, and $K_4$ block constructions}
\label{sec:rhost-block-barrier}

This section concerns a construction method, not a lower bound for arbitrary pancyclic subgraphs.
The method supplies its short cycle lengths inside separate cycle, theta, and subdivided-$K_4$ blocks.
Its cost is the sum of the cycle ranks of those blocks.
The following obstruction shows that this method cannot attain coefficient one under the stated density constraints.

\subsection{An elementary obstruction to the proposed balanced mixture}

\begin{lemma}\label{lem:rhost-excess-length}
Let $d>0$.
For a subdivision of $K_4$ with total arm length $E\ge2d$, suppose every arm has length at most $E-d$.
Then
\[
 \sum_C(|C|-d)_+\le4E-6d.
\]
For a theta graph of total length $E\ge d$,
\[
 \sum_C(|C|-d)_+\le2E-2d.
\]
Here $x_+=\max\{x,0\}$, and each sum is over all cycles of the stated graph.
\end{lemma}

\begin{proof}
Each arm of $K_4$ belongs to four of its seven cycles.
It suffices to prove $\sum_C\min\{|C|,d\}\ge6d$.
The left side is a concave function of the six nonnegative arm lengths.
For fixed $E\ge2d$, minimise it on the polytope
\[
 \sum_iw_i=E,\qquad 0\le w_i\le E-d.
\]
Its vertices have two nonzero coordinates, equal to $E-d$ and $d$.
When $E=2d$, both coordinates equal $d$.
Any two distinct edges of $K_4$ occur together in exactly two cycles.
They each occur in four cycles, so their union meets exactly six cycles.
At each polytope vertex, the sum of the truncated cycle lengths is therefore $6d$.
Concavity proves the bound throughout the polytope.

For a theta graph, every arm belongs to two cycles.
On the simplex of three nonnegative arm lengths summing to $E\ge d$,
concavity gives $\sum_C\min\{|C|,d\}\ge2d$.
Indeed, at every simplex vertex, the two nonzero cycle lengths both equal $E$.
Subtract these truncated sums from the full cycle-length sums.
\end{proof}

Consider the proposed mixture of $N+O(1)$ subdivided-$K_4$ blocks and $N+O(1)$ theta blocks.
Let $d=5N$, and suppose their arm lengths are $O(N)$.
Require the $K_4$ blocks to have total length at least $10N$ and every rank-two core at least $5N$.
Require each theta block to have total length at least $5N$.
Suppose their cycles cover all lengths $3,\ldots,10N$.
There are only $10N+O(1)$ cycle occurrences, so all but $O(1)$ must supply distinct required lengths.
Consequently, their total cycle-length sum is $50N^2+O(N)$.
Lemma~\ref{lem:rhost-excess-length} bounds their total excess above $5N$ by
\[
 50N^2-6(5N)N-2(5N)N+O(N)=10N^2+O(N).
\]
But the required lengths have excess sum
\[
 \sum_{\ell=5N+1}^{10N}(\ell-5N)
 =\frac{25}{2}N^2+O(N),
\]
a contradiction for large $N$.
This excludes arbitrary choices and parameterisations of these blocks.
It is not restricted to affine families or to a particular parity assignment.

\subsection{A quantitative barrier for arbitrary mixtures}

Define a nonnegative weight function on $[0,\infty)$ by
\[
 w(x)=
 \begin{cases}
 1,&0\le x\le1/3,\\
 4/3-x,&1/3\le x\le1,\\
 1/2-x/6,&1\le x\le2,\\
 0,&x>2.
 \end{cases}
\]
Its two formulas agree at $1/3$ and $1$.
At $2$, use the third formula.

\begin{lemma}[Finite rational certificate]\label{lem:rhost-block-weight}
Let $Q$ be a subdivision of $K_4$ with total length at least $2d$ and every rank-two core of length at least $d$.
Then
\[
 \sum_{C\subseteq Q}w(|C|/d)\le3.
\]
For a theta graph $Q$ of total length at least $d$, the corresponding sum is at most two.
For a single-cycle block, it is at most one.
\end{lemma}

\begin{proof}
The single-cycle assertion follows from $0\le w\le1$.
The other assertions have finite rational certificates, described explicitly below.
Scale the arm lengths by $d$ and write them as a nonnegative vector $x$.
For $K_4$, the metric inequalities are
\[
 \sum_{i=1}^6x_i\ge2,\qquad
 \sum_{i\ne j}x_i\ge1\quad(1\le j\le6).
\]
For a theta graph, the metric inequality is $x_1+x_2+x_3\ge1$.

Assign each cycle length to one of the four intervals in the definition of $w$.
There are $4^7=16384$ assignments for $K_4$, and $4^3=64$ for a theta graph.
Each assignment gives rational linear inequalities $Ax\le b$ and a linear objective $c\cdot x+c_0$.
The objective is the sum of cycle weights on that region.
For every assignment, the supplied certificate gives a rational vector $y$ satisfying
\[
 y\ge0,\qquad A^{\mathsf T}y\ge c,\qquad
 b\cdot y+c_0\le r,
\]
where $r=3$ for $K_4$ and $r=2$ for a theta graph.
Since $x\ge0$, these inequalities imply
\[
 c\cdot x+c_0\le y^{\mathsf T}Ax+c_0\le b\cdot y+c_0\le r.
\]
Every permitted metric lies in a listed region.
At the discontinuity $x=2$, use the region with the nonzero boundary weight.
Thus the certificates prove the claims for all real nonnegative arm lengths.
\end{proof}

The file \nolinkurl{random-host-rank3-barrier-certificates.json} contains the rational multipliers.
The standard-library verifier \nolinkurl{check_random_host_rank3_barrier.py} checks every inequality exactly.
Its optional generation mode uses a linear-programming solver only to find candidate multipliers.
Every candidate is converted to rational numbers and checked exactly before it is accepted.
The proof does not use floating-point solver success as a certificate.

\begin{proposition}[A barrier of $37/36$]\label{prop:rhost-block-barrier}
Suppose a collection of single-cycle, theta, and subdivided-$K_4$ blocks supplies all lengths $3,\ldots,\lfloor2d\rfloor$ internally.
Suppose each theta block has total length at least $d$.
Suppose each $K_4$ block has total length at least $2d$ and each of its rank-two cores has length at least $d$.
Then the sum of the block cycle ranks is at least
\[
 \frac{37}{36}d-O(1).
\]
\end{proposition}

\begin{proof}
The weight function is nonnegative.
Select one block cycle of each required length and apply Lemma~\ref{lem:rhost-block-weight} to all blocks.
Their rank sum is at least
\[
 \sum_{\ell=3}^{\lfloor2d\rfloor}w(\ell/d)
 =d\int_0^2w(x)\,dx+O(1).
\]
The function has bounded variation, so the Riemann-sum error is bounded independently of $d$.
Finally,
\[
 \int_0^2w(x)\,dx
 =\frac13+\frac49+\frac14=\frac{37}{36}.
\]
\end{proof}

The obstruction permits arbitrary numbers, metrics, and parameterisations of the three block types.
It also permits duplicate and unused cycle lengths.
It does not cover constructions using larger kernels or short cycles that traverse several blocks.
Thus it supports the search for wheel blocks without excluding coefficient one for the original random-host problem.
